% 共享输入沿同一干线分叉；两个分量求和，不表示两套独立模型。
\begin{tikzpicture}[
  x=1mm,y=1mm,font=\figurefont\small,
  box/.style={draw=BookInk,fill=BookPaper,line width=0.8pt,
    align=center,inner sep=1.5mm,minimum height=14mm},
  flow/.style={-{Latex[length=2.2mm,width=1.5mm]},
    draw=BookInk,line width=1pt,rounded corners=1.2mm},
  wire/.style={draw=BookInk,line width=1pt,rounded corners=1.2mm}
]
  \path[use as bounding box] (0,-3) rectangle (112,73);
  \node[box,text width=94mm,fill=BookBlue!5] (input) at (56,63)
    {上下文、日志行动与结果\\目标策略、日志倾向、同一个结果模型};
  \node[box,text width=45mm] (prediction) at (26,34)
    {目标策略下的预测\\$\sum_a\pi(a\mid\bm x)\widehat m(\bm x,a)$};
  \node[box,text width=45mm] (residual) at (86,34)
    {实际行动的加权残差\\
     $\dfrac{\pi(a\mid\bm x)}{\widehat\mu(a\mid\bm x)}
       [y-\widehat m(\bm x,a)]$};
  \node[box,text width=78mm,fill=BookTeal!6] (output) at (56,6)
    {两项相加得到$\psi_n$，再按评价单位求平均\\
      $\widehat V_{\mathrm{DR}}=N^{-1}\sum_n\psi_n$};
  \draw[wire] (input.south) -- (56,49);
  \draw[flow] (56,49) -| (prediction.north);
  \draw[flow] (56,49) -| (residual.north);
  \draw[wire] (prediction.south) -- (26,20) -- (56,20);
  \draw[wire] (residual.south) -- (86,20) -- (56,20);
  \draw[flow] (56,20) -- (output.north);
\end{tikzpicture}
